\subsection{The Vigil and the Holy Night}\label{sec:vigil}

The morning of 24 December still belongs to Advent. Its Mass reads Nathan's
promise to David, the same reading the Fourth Sunday has in Year B. The
\emph{Ordo lectionum Missae} (1981), the Latin edition of the Lectionary, in a
note before the Masses of 17 December, allows a reading passed over or repeated
in these days to give way to another; so when the Fourth Sunday of Year B falls
on 20 December, Isaiah's sign of the Virgin, passed over that day, may be read on
the morning of the 24th in place of Nathan. \emph{Finding the Day}, in
Appendix~\ref{sec:daily}, works out each case. The Gospel is the song Zachary
sang when his tongue was loosed, the song the Church sings every morning at Morning Prayer. It
ends with the dawn:

\begin{quote}
Through the bowels of the mercy of our God, in which the Orient from on high
hath visited us: To enlighten them that sit in darkness and in the shadow of
death: to direct our feet into the way of peace.\par
{\smallskip\noindent\small ---Luke 1:78--79\par}
\end{quote}

\noindent The Latin is \latin{oriens ex alto}, the rising from on high. When 24
December is a weekday, the verse before the Gospel that morning calls on Christ
by the same name, the Dawn, and that evening the Church will greet the light
that has risen. When it is the Fourth Sunday, the morning belongs to the
Sunday; the Sunday and Christmas are each days of obligation, and the faithful
follow their parish's schedule of Masses for both.

Advent ends before the first Evening Prayer of Christmas (Universal Norms,
n.~40). From then on it is Christmas, and the Church keeps the Nativity of the Lord with four Masses, each
with its own readings. At the Vigil Mass, said on the evening of the 24th,
Isaiah sings of Sion's marriage: ``And the bridegroom shall rejoice over the
bride, and thy God shall rejoice over thee''; St Paul preaches at Antioch of
Pisidia on the seed of David; and the Gospel is the genealogy and the dream of
Joseph from St Matthew. At the Mass during the Night the prophet sees light in
the darkness:

\begin{quote}
The people that walked in darkness, have seen a great light: to them that dwelt
in the region of the shadow of death, light is risen.\par
{\smallskip\noindent\small ---Isaiah 9:2 (9:1 in the Lectionary)\par}
\end{quote}

\noindent St Paul tells Titus that ``the grace of God our Saviour hath appeared
to all men,'' and St Luke tells of the census, the journey to Bethlehem, the
manger and the shepherds. At the Mass at Dawn the shepherds go over to
Bethlehem, and at the Mass during the Day the Church reads the prologue of St
John: the Word was made flesh.

At the Mass during the Night the angels sing:

\begin{quote}
And suddenly there was with the angel a multitude of the heavenly army, praising
God and saying: Glory to God in the highest: and on earth peace to men of good
will.\par
{\smallskip\noindent\small ---Luke 2:13--14\par}
\end{quote}

\noindent Their song is the first line of the \latin{Gloria in excelsis}, and the
Church has kept it back. On every Sunday of Advent the Missal directs that the
Gloria is not said, though the Creed is; the General Instruction gives the
Gloria to the Sundays of the year outside Advent and Lent (n.~53). The
commentary on the calendar gives the reason: the hymn
is not silenced in Advent as it is in Lent, for penance, but so that the
angels' song may sound new on Christmas night. William Durandus, the
thirteenth-century bishop of Mende who wrote at length on the meaning of the
liturgy, had said much the same:
\latin{Tempore aduentus tacetur, Gloria in excelsis, quod pertinet ad
natiuitatem, vt in nocte natiuitatis deuotius \& auidius decantetur}. In Advent
the Gloria, which belongs to the Nativity, is kept silent, so that on the night
of the Nativity it may be sung with more devotion and more longing.

On Christmas night the Church sings the \latin{Gloria in excelsis} again. It
begins with the angels' song, which the Latin Gospel of St Luke gives as
\latin{Gloria in altissimis Deo, et in terra pax hominibus bonæ voluntatis}.
